\project{14}{The IKS5A1: low and high g at once, wide pressure}{NUCLEO-H7A3ZI-Q with X-NUCLEO-IKS5A1}{Simultaneous acceleration ranges, dual full-scale barometer}

\begin{keyfacts}
\item[Board] NUCLEO-H7A3ZI-Q with X-NUCLEO-IKS5A1 on the Arduino header, and no second shield
\item[Peripherals] I2C1 for every sensor, two sensor interrupt lines on free Zio pins, TIM2 as the single timebase, USART3 for the console
\item[Toolchain] arm-none-eabi-gcc with CMake, the vendor's register-level sensor drivers vendored under their own licence
\item[Operating system] Bare metal
\item[Difficulty] 4 of 5
\item[Effort] 4 evenings of about four hours
\item[Deliverable] A shock and altitude logger that records a fine-resolution acceleration window and a wide-range pressure trace around an event the fine channel cannot see, because the coarse channel and the fine channel are both running the whole time
\end{keyfacts}

\subsection*{Why this project}

Every other sensor chapter in this book picks a range. This one refuses to. The
part that gives this shield its character reports a low-range and a high-range
acceleration measurement at the same time, from the same die, on the same
timebase, with no switching between them. That single property changes what a
logger can promise. A logger with one channel either resolves the quiet
behaviour and saturates on the event, or survives the event and cannot see the
quiet behaviour. A logger with both channels running does neither, and the
record it writes is one record rather than two recordings that have to be
aligned afterwards.

The arithmetic is worth doing before any code. Take a sixteen-bit two's
complement output, which is what these parts report, and assume the full scale
is symmetric. At a full scale of 16 g there are 32768 counts on each side, so
one count is about 0.49 mg. At a full scale of 256 g one count is about 7.8 mg.
A vibration signature at 10 mg is a little over twenty counts on the fine
channel and barely more than one count on the coarse one. That is the whole
argument for simultaneous ranges in two numbers, and it is the reason the
chapter exists. Confirm both full-scale lists and the output word format in the
part's datasheet before you trust the constants; the shape of the argument does
not change, but the constants might.

The second property is the barometer, which has two full scales rather than one:
a narrow scale for atmospheric work, where the finer sensitivity earns you
altitude resolution, and a wide scale that keeps reporting where the narrow one
stops. Pressure is the cheapest altitude estimate available and the only one on
this bench that costs no radio, and a sealed enclosure or a short water column
leaves the atmospheric range entirely. The chapter treats the range selection as
a runtime decision with a cost attached to it rather than as a configuration
constant set once at bring-up.

\begin{note}{What is actually fitted to this shield}
The X-NUCLEO-IKS5A1 carries the ISM6HG256X, an inertial measurement unit with
simultaneous low-g and high-g acceleration detection; the ISM330IS, an inertial
unit with the in-sensor programmable core; the IIS2DULPX, an ultra-low-power
accelerometer; the IIS2MDC, a magnetometer; and the ILPS22QS, a barometer with
two selectable full scales of 1260 hPa and 4060 hPa. Four independent sources
agree on that list, the vendor's own product page among them, and UM3550 is the
reference. Two parts named in older bench notes are not on this board. The
LSM6DSV32X is a consumer part on an industrial shield and reaches these boards
only through an adapter kit. The LSM6DSO16IS is the other shield's
programmable-core part, and since the two are siblings with the same core and
different qualification, it is the likely source of the confusion. Chapter 16
uses that other part deliberately.
\end{note}

Confirming that on the silkscreen takes one minute and is Step 1. Read the part
markings, run an address scan, and compare the result against UM3550. Then stop
looking at the shield and spend the rest of the chapter on the measurement,
which is where the engineering is.

\subsection*{Prior art and what to reuse}

\begin{table}[H]\centering\small
\begin{tabular}{p{34mm}p{46mm}p{38mm}p{20mm}}
\toprule
Source & What it gives & What it does not & Licence \\
\midrule
The vendor's register-level driver collection & A complete register map and a typed function per setting for every part on this shield, behind a two-function porting shim & No board layer, no bus arbitration, no ordering between two sensors, no build system, and no policy about what to record & BSD-3-Clause \\
The shield user manual UM3550 & The fitted parts, the addresses, the jumper options and which interrupt line reaches which header pin & Nothing about how to run two acceleration ranges as one measurement & Vendor document \\
The other operating system's in-tree shield definitions & The clearest published statement of how these shields route a bus and an interrupt, for the sibling shield & Whether a definition exists in that tree for this shield is checked at bring-up rather than assumed & Apache-2.0 \\
The vendor's expansion package & Working examples for the same silicon, and a sensor-fusion library & Its fusion component ships as pre-compiled objects under a licence tied to this vendor's silicon, so it is read and never vendored & Mixed, flagged \\
Chapter 13 of this book & The FIFO, the watermark interrupt and the three bus topologies, already built and measured & Nothing about a second range or a second full scale, which is this chapter & This book \\
\bottomrule
\end{tabular}
\caption{Prior art for chapter 14. The driver collection is permissive and is vendored with its copyright headers intact; the fusion middleware is not, and the difference is stated in the repository README rather than discovered by a reader later.}
\end{table}

What is left to write is the part that no source supplies: the policy. The
drivers will set any register you name. Nothing in them decides which channel
arms the recording, how long a pre-trigger window is, when the barometer changes
scale, what a record contains, or what happens when two interrupt lines assert
within the same microsecond. That policy is the chapter, and it is also the
thing a reader can carry to a different sensor next year.

\subsection*{Parts from the inventory}

\begin{table}[H]\centering\small
\begin{tabular}{p{40mm}p{66mm}p{40mm}}
\toprule
Part & Role & Interface \\
\midrule
NUCLEO-H7A3ZI-Q & Bus master, event engine, record writer, console & Micro USB to the host \\
X-NUCLEO-IKS5A1 & The five sensors, fitted to the Arduino header. The only shield on the board & I2C1 plus two interrupt lines \\
A USB data cable & Power, programming and console & Micro USB \\
nRF-PPK2 & Current, and its digital inputs used as the bench's logic recorder for the two interrupt lines & IDD jumper and two digital inputs \\
Host PC & Build, flash, terminal and the record decoder & Windows with the Linux subsystem, or a Raspberry Pi \\
\bottomrule
\end{tabular}
\caption{Inventory items used in chapter 14. Nothing is bought and nothing is soldered. One shield is on the board at a time, because two shields collide on addresses and on the 3V3 budget.}
\end{table}

\begin{note}{One shield at a time}
The three shields on this bench are not stackable in practice even though they
are stackable mechanically. Several of the parts across them answer on the same
addresses, and an address that two devices acknowledge is worse than an address
that nobody answers, because the bus reports success and the bytes are a
blend of two devices. The 3V3 rail on the board is the second reason: it is
sized for the board and one shield. This rule comes from the kit labs in the
sibling volume and is treated as authoritative here rather than re-derived.
\end{note}

\subsection*{System architecture}

\diagram{c14_arch}{One bus, five sensors, two interrupt lines and one timebase. The two acceleration channels come from the same part and therefore share a clock, which is the property the whole chapter is built on.}

The shape is conventional and the interesting part is what is not in it. There
is no fusion library, no orientation estimate and no filtering on the board
beyond a scale conversion. The firmware's job is to decide what to keep and to
keep it with an honest timestamp. Everything else can be done on the host from
the record, and anything done on the board that could have been done on the host
costs energy in the place where energy is scarce.

\subsection*{Peripheral configuration}

\begin{table}[H]\centering\small
\begin{tabular}{p{22mm}p{26mm}p{24mm}p{30mm}p{30mm}}
\toprule
Peripheral & Mode & Clock source & Pins and function & Interrupt and transfers \\
\midrule
I2C1 & Fast mode, 400 kHz & Peripheral bus, source selected in RCC & D14 and D15 on the Arduino header, which are PB9 and PB8 by the Nucleo-144 convention. Confirm in MB1363 & Interrupt-driven transfers. No DMA in this chapter \\
GPIO input & Rising edge, external interrupt & Peripheral bus & Low-g watermark line from the ISM6HG256X & EXTI, pre-emption priority 6 \\
GPIO input & Rising edge, external interrupt & Peripheral bus & High-g event line from the same part & EXTI, pre-emption priority 5 \\
TIM2 & Free running, 32-bit, 1 MHz tick & Peripheral bus & None & Read by both handlers. No interrupt \\
GPIO output & Push-pull marker pins & Peripheral bus & Two free Zio pins into the PPK2 digital inputs & None \\
USART3 & Asynchronous, 115200 8N1 & Peripheral bus & The console pins of chapter 1 & Polled \\
\bottomrule
\end{tabular}
\caption{Peripheral configuration. The high-g line is given the higher pre-emption priority because it is rare and time critical, while the watermark line is frequent and tolerant of a few microseconds of latency. The Arduino header mapping is convention until it is read in MB1363.}
\end{table}

The bus runs at 400 kHz rather than 100 kHz because the FIFO drain is the
chapter's throughput constraint and the difference is a factor of four in the
time the handler occupies. Three things decide whether that rate is actually
available. The pull-up resistors on the shield set the rise time, and they are
the shield's and not yours, so the practical ceiling is a property of the board
you are holding. The peripheral's timing register has to be programmed from the
peripheral clock that RCC is actually delivering, which on this part is not the
same branch as on its better-known sibling and is read in RM0455. And any device
on the bus may stretch the clock, so a transfer that completes in a computed
time on paper does not necessarily complete in that time on the bench. The
chapter computes the expected drain time, then measures it, and the difference
between the two is a result rather than an inconvenience.

TIM2 is the single timebase for both channels and for the barometer, and it is
32 bits wide, which at a 1 MHz tick wraps in about 71.6 minutes. Every record
carries the raw tick and the host does the conversion. Timestamping in the
handler rather than in the main loop is the whole point: a timestamp taken after
the record has been assembled measures when the firmware got around to it, not
when the sensor asserted.

\subsection*{Wiring}

\diagram{c14_wiring}{The shield sits on the Arduino header and nothing is wired by hand except the two marker pins into the power instrument. The interrupt lines reach header pins through the shield itself; which pins they reach is read from UM3550 rather than assumed.}

Two safety points before the shield goes on. Everything here is 3.3 V logic and
no 5 V signal reaches a GPIO on this board. The PPK2 measures up to 1 A, which
is far above anything this pair draws, but its digital inputs have their own
threshold and are read as logic only; they are not an analogue instrument. There
is no soldering iron on this bench, so any option on this shield that needs a
solder bridge moved is out of scope and the chapter says which configuration it
is using instead.

\subsection*{Memory and timing budget}

\diagram{c14_mem}{The record and the ring. The pre-trigger ring holds raw low-g samples so that the window before an event survives the event; the record written to the log holds both channels and the pressure trace with one timebase across all three.}

\begin{table}[H]\centering\small
\begin{tabular}{p{44mm}p{28mm}p{28mm}p{30mm}}
\toprule
Quantity & Budget & Measured & Margin \\
\midrule
Pre-trigger ring, 2048 low-g samples & 12 kB of SRAM & not measured & not measured \\
One event record, 500 ms window & about 5 kB & not measured & not measured \\
Firmware image including five drivers & 48 kB of flash & not measured & not measured \\
High-g line to first stored sample & 200 \textmu{}s & not measured & not measured \\
Barometer full-scale change to first trustworthy sample & 50 ms & not measured & not measured \\
Charge per stored event & not budgeted & not measured & not measured \\
\bottomrule
\end{tabular}
\caption{The budget table. The ring figure is arithmetic: 2048 samples of six bytes is 12288 bytes. The window figure is arithmetic too: 500 ms at 1667 Hz is 834 samples of six bytes, about 5 kB. Everything in the measured column stays blank until the DWT cycle counter or the PPK2 has produced it.}
\end{table}

The interesting budget line is the last timing row. Changing the barometer's
full scale is not free: the sensitivity constant changes, and the first
conversions after the change are taken under the new configuration but not
necessarily settled. Treating the change as instantaneous produces a pressure
trace with a step in it that looks exactly like a real event, which is the worst
kind of error because it is plausible. The chapter measures the settling
interval with the data-ready line into a PPK2 digital input and then encodes the
answer as a blanking period in the firmware.

\subsection*{Firmware design (UML)}

\diagram{c14_uml}{The event sequence. The low-g channel fills a ring continuously; the high-g line is what decides that a stretch of that ring is worth keeping. The barometer is asked for a wide-scale trace only inside the event window.}

The design is an armed recorder, not a sampler. In the quiet state the low-g
channel streams into a circular buffer through the FIFO and the watermark
interrupt built in Chapter 13, and nothing is stored. The high-g line asserts on
a threshold crossing configured in the sensor rather than tested in firmware,
which matters because a threshold tested in firmware needs the samples to have
arrived first, and by then the leading edge of the event is already inside the
ring rather than ahead of it. On assertion the firmware marks the ring position,
starts a post-trigger countdown, changes the barometer to its wide scale, waits
out the blanking period, and then assembles one record from the ring, the high-g
samples and the pressure trace.

Two ordering rules keep it honest. First, the ring is never rewound: the
pre-trigger window is whatever the ring already held, and if the ring is shorter
than the requested window the record says so in a flag rather than silently
returning a short window. Second, a second high-g assertion inside an open
window extends the window rather than starting a second record, because two
overlapping records of one physical event are harder to interpret than one long
one.

\subsection*{Data flow (ASCII)}

\begin{asciiart}
  X-NUCLEO-IKS5A1                       STM32H7A3                         host
  +---------------------------+         +-----------------------------+   +-------------+
  | ISM6HG256X                |  I2C1   | i2c1_xfer (interrupt)       |   | decoder.py  |
  |   low-g path  -> FIFO     +-------->|   |                         |   |   |         |
  |   high-g path -> threshold|         |   v                         |   |   v         |
  |                           |  INT1   | ring[2048] 6 B per sample   |   | CSV + plot  |
  |   INT1 watermark  --------+-------->|   |                         |   |             |
  |   INT2 high-g     --------+-------->|   +--> arm: mark, countdown  |   |             |
  +---------------------------+  INT2   |        |                    |   |             |
  | ILPS22QS 1260 / 4060 hPa  |  I2C1   |        v                    |   |             |
  |   data-ready      --------+-------->| record assembler            |   |             |
  +---------------------------+         |   |                         |   |             |
                                        |   v                         |USART3           |
  PPK2 digital in D0 <-- marker A       | framed record ---------------+-->|             |
  PPK2 digital in D1 <-- marker B       +-----------------------------+   +-------------+
\end{asciiart}

\subsection*{Repository layout}

\begin{plaincode}
nucleo-h7a3-iks5a1-shocklog/
  CMakeLists.txt
  cmake/arm-none-eabi.cmake
  third_party/stmems/           # vendored, BSD-3-Clause, headers untouched
    ism6hg256x_reg.[ch]
    ilps22qs_reg.[ch]
    iis2mdc_reg.[ch]
    LICENSE                     # copied verbatim, not summarised
  src/bus_i2c1.c                # the two-function porting shim
  src/shock_cfg.c               # ranges, rates, thresholds, in one place
  src/ring.c                    # pre-trigger ring, from chapter 2
  src/recorder.c                # the policy: arm, extend, assemble
  src/baro_range.c              # full-scale selection and its blanking period
  src/timebase.c                # TIM2 at 1 MHz, read from both handlers
  src/main.c
  tools/decode_record.py        # the host twin of the record format
  tools/scan_i2c.py             # address scan, run once at bring-up
  docs/PROVENANCE.md            # what is vendored, under which licence
  README.md
\end{plaincode}

\subsection*{Steps}

\begin{steps}
\item \textbf{Confirm the shield before writing a line of driver code.} Read the
part markings on the silkscreen and compare them against UM3550. You are looking
for ISM6HG256X, ISM330IS, IIS2DULPX, IIS2MDC and ILPS22QS. If you find an
LSM6DSV32X or an LSM6DSO16IS you are holding a different board, and every
address and register in this chapter is wrong for it. Then confirm electrically
with a scan, because a marking you misread and an address that answers are two
independent pieces of evidence.
\begin{shellcode}
python tools/scan_i2c.py --port /dev/ttyACM0
# prints every address that acknowledges, one line each
\end{shellcode}
Write the addresses the scan returns into a header with a comment naming UM3550
as the source. This book does not print sensor addresses it has not read on the
bench, and neither should your repository.

\item \textbf{Vendor the register-level drivers and write the shim.} The
collection is BSD-3-Clause, so it can be copied into the repository with its
copyright headers intact. Record that in \texttt{docs/PROVENANCE.md} on the same
day you copy it, not later. Each driver needs exactly two functions from you.
\begin{ccode}
/* The porting shim. One instance per sensor, differing only in address. */
typedef struct { uint16_t addr; } bus_ctx_t;

int32_t stmems_write(void *h, uint8_t reg, const uint8_t *buf, uint16_t n)
{
    const bus_ctx_t *c = (const bus_ctx_t *) h;
    return i2c1_write_reg(c->addr, reg, buf, n) ? 0 : -1;
}

int32_t stmems_read(void *h, uint8_t reg, uint8_t *buf, uint16_t n)
{
    const bus_ctx_t *c = (const bus_ctx_t *) h;
    return i2c1_read_reg(c->addr, reg, buf, n) ? 0 : -1;
}
\end{ccode}
The driver's own configuration function names are taken from its header when you
compile against it. This book does not print an API it has not compiled, which
is why the configuration below goes through a wrapper of your own.

\item \textbf{Configure both acceleration channels and prove they are both
live.} This is the step the chapter exists for. One structure holds both, and
neither is a mode you switch into.
\begin{ccode}
typedef struct {
    uint16_t lo_odr_hz;    /* fine channel rate, e.g. 1667 */
    uint16_t lo_fs_g;      /* fine channel full scale, e.g. 16 */
    uint16_t hi_fs_g;      /* coarse channel full scale, from the datasheet */
    uint16_t hi_thresh_mg; /* threshold evaluated in the sensor, not here */
    uint16_t fifo_wtm;     /* watermark in samples, see chapter 13 */
} shock_cfg_t;

/* Acceptance for this step: with the board held still, the fine channel
   reports about 1 g on one axis with counts of noise, and the coarse
   channel reports the same 1 g with far fewer counts of resolution.
   Both at the same time, from one part, on one timestamp. */
\end{ccode}
If only one channel ever reports, you have configured a range rather than two
channels, and the rest of the chapter does not work. Fix it here.

\item \textbf{Bring up the FIFO and the watermark line.} The mechanism is
Chapter 13's and is not repeated: the sensor fills its own FIFO, asserts a line
at the watermark, and the handler drains a known number of samples rather than
polling a status bit. What is new here is where the samples go, which is a
circular pre-trigger ring rather than a consumer.
\begin{ccode}
void EXTI_wtm_handler(void)
{
    const uint32_t t = timebase_ticks();      /* TIM2, before any bus work */
    uint8_t raw[FIFO_WTM * 6];
    (void) i2c1_read_fifo(ACC_ADDR, raw, sizeof raw);
    ring_push_block(&pretrig, raw, FIFO_WTM, t);
    marker_pulse(MARKER_A);                   /* PPK2 digital input D0 */
}
\end{ccode}

\item \textbf{Arm the recorder on the high-g line.} The threshold lives in the
sensor. The handler does almost nothing, because the handler runs at the moment
that matters and everything it does is latency in the record.
\begin{ccode}
volatile uint32_t g_event_tick;
volatile uint8_t  g_event_pending;

void EXTI_highg_handler(void)
{
    g_event_tick = timebase_ticks();
    g_event_pending = 1;                      /* the main loop assembles */
    marker_pulse(MARKER_B);                   /* PPK2 digital input D1 */
}
\end{ccode}
Measure the interval from the marker pulse to the first byte of the assembled
record with the DWT cycle counter, and cross-check the two marker pulses on the
PPK2 digital inputs, which are time-aligned with its current trace. That
cross-check is what this bench has instead of a logic analyser, and it is enough
for intervals of this length.

\item \textbf{Change the barometer's full scale and find out what that costs.}
Select the narrow 1260 hPa scale for the quiet state, where the finer
sensitivity is what gives altitude any resolution, and the wide 4060 hPa scale
inside an event window. Then measure the cost of the change rather than assuming
it is free.
\begin{ccode}
/* Blanking: samples taken inside this window are tagged, not discarded,
   so the host can see what the firmware chose to distrust. */
#define BARO_BLANK_US  50000u      /* placeholder until measured */

void baro_set_range(baro_range_t r)
{
    baro_apply_range(r);                       /* wrapper over the driver */
    g_baro_valid_from = timebase_ticks() + BARO_BLANK_US;
}
\end{ccode}
To measure it: hold the pressure steady, change the scale, and record every
sample with its timestamp. The interval from the change to the point where the
reading stops moving is the blanking period. Replace the placeholder with the
measured figure and say in the README which instrument produced it.

\item \textbf{Define the record and write its host twin at the same time.} One
record, one timebase, both channels, explicit flags for everything the firmware
distrusted. Write the decoder in the same sitting, because a format with no
second implementation is a format with no specification.
\begin{ccode}
typedef struct __attribute__((packed)) {
    uint32_t tick_event;      /* TIM2 ticks at the high-g assertion */
    uint16_t n_pre;           /* fine-channel samples before the event */
    uint16_t n_post;          /* fine-channel samples after it */
    uint16_t n_baro;          /* pressure samples in the window */
    uint8_t  lo_fs_g;
    uint8_t  hi_fs_g;
    uint8_t  baro_range;      /* 0 = 1260 hPa, 1 = 4060 hPa */
    uint8_t  flags;           /* bit0 short window, bit1 extended, bit2 baro blanked */
} record_hdr_t;
\end{ccode}

\item \textbf{Produce a controlled event and check the record against it.} No
instrument on this bench generates a calibrated shock, so do not pretend to one.
A repeatable event is enough: let the board and shield, still attached to their
cable, settle onto a foam pad from a fixed low height, ten times. The record
should show a quiet fine channel, a coarse-channel excursion at the moment the
marker pin pulsed, and a pressure trace that is unremarkable, which is itself
the result. Then repeat with the assembly inside a sealed bag that you squeeze,
which moves the pressure and leaves the acceleration alone, and check that the
two channels are independent in the record.

\item \textbf{Decode and plot on the host.} The decoder reads the framed record
from the console, converts counts to engineering units using the scale codes in
the header rather than a constant compiled into the script, and derives altitude
from pressure with the standard relation. Keeping the scale in the record is
what lets a record recorded under one configuration be read years later.
\begin{pycode}
# altitude from pressure, the standard barometric relation
h_m = 44330.0 * (1.0 - (p_hpa / p0_hpa) ** 0.1903)
\end{pycode}

\item \textbf{Find the rate at which the bus stops keeping up.} The fine channel
rate is the one number in the configuration that has a hard ceiling, and the
ceiling is not in the sensor. Compute the expected drain time first, so that you
know what you are looking for.
\begin{shellcode}
# 128 samples of 6 bytes is 768 bytes, plus address and register overhead.
# At 400 kHz each byte costs 9 bit times, so 768 x 9 / 400000 = 17.3 ms.
# Computed from the bus rate. Not measured.
\end{shellcode}
Then raise the rate in steps, and at each step record two numbers with the DWT
cycle counter: the time spent inside the watermark handler, and the fraction of
the sample period that represents. When the handler occupies more than about a
quarter of the period, the design has stopped being a logger and has become a
bus benchmark, and the answer is either a larger watermark, a faster bus, or the
transfer engine of Chapter 4. Record where that point is, because it is the
number a reader of your repository will want and nobody has published it for
this shield.
\end{steps}

\subsection*{Build, flash and debug}

\diagram{c14_timing}{The two interrupt lines and the window they define. The pre-trigger portion is already in the ring when the high-g line asserts, which is the only way to record the leading edge of an event the fine channel cannot detect in time.}

\begin{shellcode}
cmake -B build-fw -G Ninja -DCMAKE_TOOLCHAIN_FILE=cmake/arm-none-eabi.cmake
cmake --build build-fw -j
cp build-fw/shocklog.bin "$PROBE_DISK"/   # onto the probe disk
python tools/decode_record.py --port /dev/ttyACM0 --out events/
\end{shellcode}

Recovering a board that will not answer follows Chapter 1: the probe presents a
disk, and copying a known-good binary onto it needs no tools at all. The failure
specific to this chapter is different, and it is a silent one.

\begin{note}{When the high-g channel reports nothing}
In order of likelihood: the second channel was never enabled, so the part is
reporting one range twice; the threshold is set in the wrong units, since a
threshold expressed in milli-g and a threshold expressed in full-scale counts
differ by whatever the full scale happens to be; the interrupt line is routed to
a header pin the firmware is not watching, which UM3550 settles in a minute; the
line is configured for the wrong edge; and the pull configuration on the input
pin opposes the shield's own. Check them in that order, and check the first one
by reading both channels back with the board held still, where a working pair
reports the same acceleration at visibly different resolutions.
\end{note}

\subsection*{Verification and acceptance criteria}

\begin{itemize}
\item With the board held still, both acceleration channels report the same
value within their respective resolutions, simultaneously, from one part, with
one timestamp. This is the chapter's central claim and it is checked first.
\item The pre-trigger window in a stored record contains samples whose
timestamps are earlier than the high-g assertion tick, proving the ring held the
leading edge rather than the firmware reconstructing it.
\item A second high-g assertion inside an open window extends that window and
sets the extension flag. Two overlapping records for one physical event never
appear.
\item The barometer's blanking period is a measured number in the README with
its instrument named, not the placeholder constant, and samples inside it are
tagged rather than discarded.
\item The host decoder reads the record using only the scale codes carried in
the header, proven by decoding a record captured under a different full-scale
configuration.
\item A pressure event with no acceleration and an acceleration event with no
pressure change both produce records in which the other channel is unremarkable.
\item Every measured figure in the repository names the DWT cycle counter, TIM2
or the PPK2. Nothing is asserted from a feeling about how long something took.
\end{itemize}

\subsection*{Variants}

\begin{table}[H]\centering\small
\begin{tabular}{p{24mm}p{26mm}p{44mm}p{22mm}p{20mm}}
\toprule
Axis & Variant & What changes & Cost & Built in full in \\
\midrule
Execution model & Polled loop & Poll the status bit instead of watching the watermark line. Simpler, and it loses the leading edge of exactly the events this chapter is about & The pre-trigger window becomes a fiction & Chapter 3 \\
Execution model & DMA into the ring & The FIFO drain becomes a transfer the core does not supervise, which matters once the fine channel rate rises & Cache maintenance on the ring buffer & Chapter 4 and Chapter 19 \\
Time and safety & RTC and a low-power timer & The quiet state stops being a busy loop and becomes a wait, which is what makes a logger a logger & Wake latency added to the pre-trigger reasoning & Chapter 7 \\
Intelligence and reach & Classifier in the sensor & The ISM330IS on this same shield carries the programmable core, so the arming decision could be made in the sensor rather than by a threshold & A toolchain and a licence audit & Chapter 16 \\
Operating system & An RTOS task set & The drain, the assembler and the console become three tasks with stated priorities and a measured latency distribution & Static allocation and a scheduler to justify & Chapter 20 \\
Peripheral substitution & A different bus instance & The same five sensors read through another I2C instance, which changes the clock domain and therefore what survives a low-power mode & Pin availability while a shield is fitted & Chapter 17 \\
\bottomrule
\end{tabular}
\caption{Variants for chapter 14. Nothing in this table is built in full here; this chapter's contribution is the simultaneous-range policy, and each variant is named so a reader knows where the full treatment lives.}
\end{table}

\subsection*{Pitfalls}

\begin{itemize}
\item Trusting an older bench note about which parts are on this shield. Two
parts that appear in such notes are not fitted, and one of them is a sibling of
a part that is, which makes the mistake survive a casual check.
\item Fitting two shields at once. Addresses collide and the 3V3 budget is
shared. One shield is on the board at a time.
\item Treating the two acceleration channels as a range selection. If your
configuration code has a branch that chooses between them, you have not
configured this part, you have configured a different one.
\item Setting a threshold in the wrong units. Milli-g and full-scale counts are
not the same number and the error scales with the full scale, so it is small at
16 g and large at 256 g.
\item Timestamping after the bus transaction. The handler reads the timebase
first, before any I2C work, or the timestamp measures the bus rather than the
event.
\item Assuming the barometer's full-scale change is instantaneous. The step it
produces in an untagged trace looks exactly like a real pressure event.
\item Compiling the scale constants into the host decoder. The record carries
its own scale codes so that it can be read later, by someone else, on a
different configuration.
\item Reading register offsets from a tutorial written for a different member of
this sensor family. The sibling parts share a core and differ in qualification,
addresses and available registers.
\end{itemize}

\subsection*{Best practices applied}

\begin{itemize}
\item The prior art is named and its licence is stated before any code is
written around it, and the vendored copy keeps its copyright headers and its
licence file verbatim rather than in summary.
\item Every measurement names its instrument, and a measurement that has not
been taken is written as not measured rather than estimated.
\item Facts about the shield carry their evidence: the part list is confirmed
from four independent sources and then again on the silkscreen and on the bus,
and the addresses are read rather than printed from memory.
\item The record format has two implementations from the first day, one on the
board and one on the host, which is the cheapest specification there is.
\item What the firmware distrusts is flagged in the data rather than removed
from it, so that a later reader can disagree with the firmware.
\end{itemize}

\subsection*{Stretch goals}

\begin{itemize}
\item Add the IIS2DULPX as a third acceleration channel and compare its noise
floor against the fine channel of the main part at the same rate. The
ultra-low-power part is not a smaller version of the other one, and the
comparison is a genuine result because nobody has published it for this shield.
\item Use the magnetometer to record orientation at the moment of the event, so
that a record says which way up the assembly was. Keep the estimate on the host
using a permissive orientation filter rather than putting fusion on the board.
\item Store records in flash rather than streaming them, which raises the
write-granularity and error-correction questions this family has and which are
on the confirm list rather than settled.
\item Drive the arming decision from the in-sensor core on the ISM330IS and
compare the latency and the charge against the threshold path built here. That
comparison is Chapter 16's subject and this shield can host half of it.
\end{itemize}

\subsection*{Roadmap and next steps}

Chapter 15 changes the sensor and keeps the discipline: a time-of-flight array
whose output is sixty-four values that nobody wants, reduced on the
microcontroller to a decision that somebody does. Chapter 16 then takes the
arming decision built here and moves it to three different places to find out
what each one costs.

For going wider, the community embedded engineering roadmap distinguishes
firmware from embedded Linux from hardware and is explicit about which to weight
at which stage, and its central claim, that projects outrank reading, is this
book's claim too. The accompanying essay on bridging the gap between the two is
the short version. For the sensor-intelligence direction this chapter points at,
the free open textbook that grew out of a university course is the closest thing
to a canonical curriculum, and the free-to-audit course on embedded machine
learning covers the same ground with exercises. Appendix H lists them with their
licences, which matters here: one widely recommended video course releases its
code under terms that make it something to recommend and never to quote.

\subsection*{Portfolio evidence}

\begin{itemize}
\item A public repository whose README opens with the architecture figure and
states in its first paragraph that this shield carries the ISM6HG256X and not
the part older notes name, with the four sources listed.
\item A plot of one event showing both acceleration channels on one time axis,
with the high-g assertion marked, and the pre-trigger portion visibly earlier
than the assertion.
\item A PPK2 capture in which the two marker pins and the current trace are
time-aligned, which is what this bench has instead of a logic analyser and is
worth saying so in the caption.
\item The provenance document listing every vendored file, its licence and the
date it was copied.
\end{itemize}

\subsection*{Sources}

Normative references:
\begin{itemize}
\item UM3550, the user manual for the X-NUCLEO-IKS5A1, for the fitted parts, the
addresses, the jumper options and the interrupt routing to header pins.
\item The datasheets of the ISM6HG256X, ISM330IS, IIS2DULPX, IIS2MDC and
ILPS22QS, for the full-scale lists, the output word format, the output data
rates and the interrupt behaviour.
\item Reference manual RM0455 for I2C1 and the external interrupt controller on
this part. Not RM0433, which documents its better-known sibling.
\item The board user manual for MB1363, for the Arduino header mapping, which
pins D14 and D15 reach, and which Zio pins stay free while a shield is fitted.
\end{itemize}

Reusable implementations:
\begin{itemize}
\item The vendor's register-level driver collection, BSD-3-Clause, covering
every sensor on this shield behind a two-function porting shim.\newline
\url{https://github.com/STMicroelectronics/STMems_Standard_C_drivers}
\item The other operating system's in-tree shield definitions, Apache-2.0, which
document the bus topologies of the sibling shield more clearly than the vendor's
own manual.\newline
\url{https://docs.zephyrproject.org/latest/boards/shields/x_nucleo_iks4a1/doc/index.html}
\item The in-sensor processor examples for the programmable core on the
ISM330IS. The repository declares no licence at its root, so files are checked
individually and nothing is vendored on faith.\newline
\url{https://github.com/STMicroelectronics/st-mems-ispu}
\item The current home of the well-known orientation filter, MIT, for the
host-side orientation stretch goal.\newline
\url{https://github.com/xioTechnologies/Fusion}
\item The Python interface to the power instrument, GPL-2.0. Install it from its
package index; do not fork it into a portfolio repository.\newline
\url{https://github.com/IRNAS/ppk2-api-python}
\end{itemize}
